\chapter{2007年江西卷（理）}

\section{单选题}

\begin{problem}
化简 \(\frac{2 + 4i}{(1 + i)^2}\) 的结果是（\quad）

\choices
    {\(2 + i\)}
    {\(-2 + i\)}
    {\(2 - i\)}
    {\(-2 - i\)}
\end{problem}

\begin{problem}
\(\lim_{x\to 1}\frac{x^3 - x^2}{x - 1}\)（\quad）

\choices
    {等于 0}
    {等于 1}
    {等于 3}
    {不存在}
\end{problem}

\begin{problem}
若 \(\tan \left(\frac{\pi}{4} - \alpha\right) = 3\)，则 \(\cot \alpha\) 等于（\quad）

\choices
    {\(-2\)}
    {\(-\frac{1}{2}\)}
    {\(\frac{1}{2}\)}
    {2}
\end{problem}

\begin{problem}
已知 \(\left(\sqrt{x} + \frac{3}{\sqrt{x}}\right)^n\) 展开式中，各项系数的和与其各项二项式系数的和之比为 64，则 \(n\) 等于（\quad）

\choices
    {4}
    {5}
    {6}
    {7}
\end{problem}

\begin{problem}
若 \(0 < x < \frac{\pi}{2}\)，则下列命题中正确的是（\quad）

\choices
    {\(\sin x < \frac{3}{2} x\)}
    {\(\sin x > \frac{3}{2} x\)}
    {\(\sin x < \frac{4}{3} x^2\)}
    {\(\sin x > \frac{4}{4}\pi x^2\)}
\end{problem}

\begin{problem}
若集合 \(M = \{0, 1, 2\}\)，\(N = \{(x, y) | x - 2y + 1 \geqslant 0 \text{ 且 } x - 2y - 1 \leqslant 0, x, y \leqslant M\}\)，则 \(N\) 中元素的个数为（\quad）

\choices
    {9}
    {6}
    {4}
    {2}
\end{problem}

\begin{problem}
如图，正方体 \(AC_{1}\) 的棱长为 1 ，过点 \(A\) 作平面 \(A_{1}BD\) 的垂线，垂足为点 \(H\). 则以下命题中，错误的是（\quad）

\choices
    {点 H 是 \( \triangle A_{1}BD \) 的重心}
    {AH 垂直平面 \( CB_{1}D_{1} \)}
    {AH 的延长线经过点 \( C_{1} \)}
    {直线 AH 和 \( BB_{1} \) 所成的角为 \( 45^{\circ} \)}
\end{problem}

\begin{problem}
四位好朋友在一次聚会上，他们按照各自的爱好选择了形状不同、内容高度相等、杯口半径相等的圆口酒杯，如图所示，盛满酒后他们的定: 先各自饮杯中酒的一半. 设剩余酒的高度从左到右依次为 \(h_1, h_2, h_3, h_4\)，则它们的大小关系正确的是（\quad）

\choices
    {\(h_2 > h_1 > h_4\)}
    {\(h_1 > h_2 > h_3\)}
    {\(h_3 > h_2 > h_4\)}
    {\(h_2 > h_4 > h_1\)}
\end{problem}

\begin{problem}
设椭圆 \(\frac{x^2}{a^2} + \frac{y^2}{b^2} = 1 (a > b > 0)\) 的离心率 \(e = \frac{1}{2}\)，右焦点为 \(F(c,0)\). 方程 \(ax^2 + bx - c = 0\) 的两个实根分别为 \(x_1\) 和 \(x_2\)，则点 \(P(x_1, x_2)\)（\quad）

\choices
    {必在圆 \(x^{2} + y^{2} = 2\) 内}
    {必在圆 \(x^{2} + y^{2} = 2\) 上}
    {必在圆 \(x^{2} + y^{2} = 2\) 外}
    {以上三种情形都有可能}
\end{problem}

\begin{problem}
将一个骰子连续抛掷三次，它落地时间上的点数依次成等差数列的概率为（\quad）

\choices
    {\( \frac{1}{9} \)}
    {\( \frac{1}{12} \)}
    {\( \frac{1}{10} \)}
    {\( \frac{1}{18} \)}
\end{problem}

\begin{problem}
设函数 \( f(x) \) 是 \( \R \) 上以 5 为周期的可导偶函数，则曲线 \( y = f(x) \) 在 \( x = 5 \) 处的切线的斜率为（\quad）

\choices
    {\(-\frac{1}{9}\)}
    {0}
    {\(\frac{1}{8}\)}
    {5}
\end{problem}

\begin{problem}
设 \( p: f(x) = \e^{x} + \ln x + 2x^{2} + mx + 1 \) 在 \((0, +\infty)\) 内单调递增，\( q: m \geqslant -5 \)，则 \( p \) 是 \( q \) 的（\quad）

\choices
    {充分不必要条件}
    {必要不充分条件}
    {充分必要条件}
    {既不充分也不必要条件}
\end{problem}

\section{填空题}

\begin{problem}
设函数 \( y = 4 + \log_2(x - 1) (x \geqslant 3) \)，则其反函数的定义域为 \fillinblank{} \fillinblank{}.
\end{problem}

\begin{problem}
已知数列 \( \{a_{n}\} \) 对任予任意 \( p, q \in N^{*} \) ，有 \( a_{p} + a_{q} = a_{p+q} \) ，若 \( a_{1} = \frac{1}{9} \) ，则 \( a_{99} = \)  \fillinblank{}.
\end{problem}

\begin{problem}
如图，在 \(\triangle ABC\) 中，点 \(O\) 是 \(BC\) 的中点，过点 \(O\) 的直线分别交直线 \(AB\)，\(AC\) 于不同的两点 \(M, N\)，若 \(\overrightarrow{AB} = m\overrightarrow{AM}, \overrightarrow{AC} = n\overrightarrow{AN}\)，则 \(m + n\) 的值为 \fillinblank{} \fillinblank{}.
\end{problem}

\begin{problem}
设有一组圆 \(C_{1}: (x - k + 1)^{2} + (y - 3k)^{2} = 2k^{4}\) (\(k \in \mathbf{N}^{*}\)). 下面四个命题:

\choices{存在一条定直线与所有的圆均相切;}{存在一条定直线与所有的圆均相交;}{存在一条定直线与所有的圆均不相交;}{所有的圆均不经过原点. 其中真命题的代号是 \fillinblank{}. (写出所有真命题的代号)}
\end{problem}

\section{解答题}

\begin{problem}
已知函数 \( f(x) = \begin{cases} cx + 1, & 0 < x < c, \\ 2^{-\frac{x}{3}} + k, & c \leqslant x < 1, \end{cases} \) 在区间 \((0,1)\) 内连续，且 \( f(x^2) = \frac{9}{5} \). (1) 求实数 \(k\) 和 \(c\) 的值; (2) 解不等式 \( f(x) > \frac{\sqrt{2}}{8} + 1 \) .
\end{problem}

\begin{problem}
如图，函数 \( y = 2\cos (\omega x + \theta) \left( x \in \R, 0 \leqslant \theta \leqslant \frac{\pi}{2} \right) \) 的图象与 \( y \) 轴交于点 \( (0, \sqrt{3}) \)，且在该点处切线的斜率为 \(-2\). (1) 求 \(\theta\) 和 \(t\) 的值; (2) 已知点 \(A\left(\frac{\pi}{2}, 0\right)\)，点 \(P\) 是该函数图象上的一点，点 \(Q(x_0, y_0)\) 是 \(PA\) 的中点，当 \(y_0 = \frac{\sqrt{3}}{2}, x_0 \in \left[\frac{\pi}{2}, \pi\right]\) 时，求 \(x_0\) 的值.
\end{problem}

\begin{problem}
某陶瓷厂准备烧制甲、乙、丙三件不同的工艺品，制作过程中必须先后经过两次烧制，当第一次烧制合格后可进入第二次烧制，两次烧制过程相互独立. 根据该厂现有的技术水平，经过第一次烧制后，甲、乙、丙三件产品合格的概率依次为0.5, 0.6, 0.4. 经过第二次烧制后，甲、乙、丙三件产品合格的概率依次为0.6, 0.5, 0.75. (1) 求第一次烧制后给有一件产品合格的概率; (2) 经过前后两次烧制后，合格工艺品的个数为 6 ，求随机变量 \( \xi \) 的期望.
\end{problem}

\begin{problem}
设动点 \(P\) 到点 \(A(-1,0)\) 和 \(B(1,0)\) 的距离分别为 \(d_{1}\) 和 \(d_{2}, \angle APB = 2\theta\). 且存在常数 \(\lambda (0 < \lambda < 1)\)，使得 \(d_{1}d_{2}\sin^{2}\theta = \lambda\). (1) 证明: 动点 \(P\) 的轨迹 \(C\) 为双曲线，并求出 \(C\) 的方程; (2) 过点 B 作直线交双曲线 C 的右支于 M、N 两点，试确定 \( \lambda \) 的范围，使 \( \overrightarrow{OM} \cdot \overrightarrow{ON} = 0 \) ，其中点 O 为坐标原点.
\end{problem}

\begin{problem}
设正整数数列 \(\{a_{n}\}\) 满足： \(a_2 = 4\) ，且对于任何 \(n\in \mathbf{N}^*\) ，有 \(2 + \frac{1}{a_{n + 1}} <\) \[ \frac {1}{a _{n}} + \frac {1}{\frac {a _{n + 1}}{1} - \frac {1}{n + 1}} < 2 + \frac {1}{a _{n}}. \] (1) 求 \(a_{1}, a_{2}\); (2) 求数列 \(\{a_{n}\}\) 的通项 \(a_{n}\).
\end{problem}

\begin{problem}
如图是一个直三棱柱（以 \(A_{1}B_{1}C_{1}\) 为底面）被一平面所截得到的几何体，截面为 \(ABC\). 已知 \(A_{1}B_{1} = B_{1}C_{1} = 1, \angle A_{1}B_{1}C_{1} = 90^{\circ}, AA_{1} = 4, BB_{1} = 2, CC_{1} = 3\). (1) 设点 O 是 AB 的中点，证明: \( OC // \) 平面 \( A_{1}B_{1}C_{1} \) ; (2) 求二面角 \(B - AC - A_{1}\) 的大小; (3) 求此几何体的体积.
\end{problem}

